\chapter{Sucesiones}\label{ch:g12:seq}

A \omterm{def:g12:seq:sequence}{sucesión} es una lista de números reales indexado por el números naturales. sucesiones
modelo de evoluciones discretas --- poblaciones contadas año tras año, saldos de un
cuenta bancaria, sucesivas aproximaciones de un número --- y sus límites son
el primer encuentro serio con el infinito. Este capítulo establece el
vocabulario, el principio de inducción y la convergencia fundamental
teoremas.

\section{Razonamiento por inducción}

\begin{theorem}[Principio de inducción]\label{thm:g12:seq:induction}
Sea $P(n)$ ser una declaración dependiendo de un entero $ n $, y deja $ n_0 \in \N $.
Si
\begin{enumerate}
  \item \emph{(caso base)} $ P(n_0)$ es verdad, y
  \item \emph{(paso inductivo)} por cada $ n \geq n_0$, $ P(n)$ implica $ P(n+1)$,
\end{enumerate}
entonces $ P(n)$ es cierto para cada $ n \geq n_0$.
\end{theorem}

\begin{proof}
Supongamos, por contradicción, que el conjunto $ A $ de números enteros $ n \geq n_0$ para
cual $ P(n)$ es falso no está vacío. Entonces $ A $ tiene un elemento más pequeño
$ m $.\footnote{Every nonempty subset of $\N $ has a smallest element; this
property of $\N $ is taken as an axiom.} Desde $ P(n_0)$ es verdad, $ m > n_0$, entonces
$ m - 1 \geq n_0$ y $ m-1 \notin A $, es decir.\ $ P(m-1)$ es verdad. el inductivo
paso aplicado a $ n = m-1$ luego \omterm{def:g10:proba:sample}{muestra} que $ P(m)$ es cierto, contradice
$ m \in A $.
\end{proof}

\begin{example}\label{ex:g12:seq:bernoulli}
Probemos \emph{La desigualdad de Bernoulli}: por cada real $ a > 0$ y cada
$ n \in \N $,
\[
(1+a)^n \geq 1 + na.
\]
\emph{Caso base.} Para $ n = 0$, ambos lados son iguales $1$.
\emph{Paso inductivo.} Suponer $(1+a)^n \geq 1+na $ para algunos $ n \in \N $. Desde
$1 + a > 0$, multiplicando ambos lados por $1+a $ preserva la desigualdad:
\[
(1+a)^{n+1} \geq (1+na)(1+a) = 1 + (n+1)a + na^2 \geq 1 + (n+1)a .
\]
Por inducción, la desigualdad es válida para todos $ n \in \N $.
\end{example}

\begin{method}[Escribir una prueba de inducción]\label{met:g12:seq:induction}
Siempre haz la declaración $ P(n)$ explícito antes de comenzar. Una prueba completa
tiene tres partes visibles: el caso base, el paso inductivo (``suponer $ P(n)$;
probamos $ P(n+1)$''), y la conclusión que invoca el principio de inducción.
El error más común es probar el paso inductivo sin utilizar nunca el
hipótesis $ P(n)$: si eso sucede, o la prueba es incorrecta o la inducción fue incorrecta
no es necesario.
\end{method}

\section{Vocabulario de sucesiones.}

\begin{definition}[Sucesión]\label{def:g12:seq:sequence}
A \emph{sucesión}\index{sucesión} es una \omterm{def:g11:func:function}{función} $ u \colon \N \to \R $ (o de
$\{n \in \N : n \geq n_0\}$ a $\R $). El \omterm{def:g10:functions:function}{imagen} de $ n $ esta escrito $ u_n $, y
la sucesión misma $(u_n)_{n\in\N}$ o simplemente $(u_n)$.
\end{definition}

A \omterm{def:g12:seq:sequence}{sucesión} puede ser definido \emph{explícitamente}, por una fórmula $ u_n = f(n)$, o por
\emph{reaparición}, por su primer término y una relación $ u_{n+1} = f(u_n)$.

\begin{definition}[monotonicidad]\label{def:g12:seq:monotonic}
A \omterm{def:g12:seq:sequence}{sucesión} $(u_n)$ es \emph{creciente}\index{creciente} si
$ u_{n+1} \geq u_n $ para todo $ n $, \emph{\omterm{def:g10:functions:variations}{decreciente}} si $ u_{n+1} \leq u_n $ para
todos $ n $, y \emph{monótono} si está aumentando o \omterm{def:g10:functions:variations}{decreciente}. Es
\emph{estrictamente} aumentando (resp.\ \omterm{def:g10:functions:variations}{decreciente}) cuando las desigualdades son
estricto.
\end{definition}

\begin{method}[Estudiar la monotonicidad de una sucesión.]\label{met:g12:seq:monotonicity}
Tres técnicas estándar:
\begin{enumerate}
  \item estudiar el signo de $ u_{n+1} - u_n $;
  \item si todos los términos son positivos, compare $\dfrac{u_{n+1}}{u_n}$ a $1$;
  \item si $ u_n = f(n)$ con $ f $ definido en $\intco{0}{+\infty}$, utiliza el
        variaciones de $ f $.
\end{enumerate}
\end{method}

\begin{definition}[sucesión acotada]\label{def:g12:seq:bounded}
A \omterm{def:g12:seq:sequence}{sucesión} $(u_n)$ es \emph{delimitado arriba}\index{delimitado arriba} si hay
existe $ M \in \R $ con $ u_n \leq M $ para todo $ n $; \emph{delimitado por debajo} si hay
existe $ m \in \R $ con $ u_n \geq m $ para todo $ n $; y \emph{encerrado} si ambos
espera.
\end{definition}

\subsection{Sucesiones aritméticas y geométricas.}

\begin{definition}[Sucesiones aritméticas y geométricas.]\label{def:g12:seq:arith-geom}
A \omterm{def:g12:seq:sequence}{sucesión} $(u_n)$ es \emph{aritmética}\index{aritmética} con común
\omterm{def:g11:prob:variance}{diferencia} $ r $ si $ u_{n+1} = u_n + r $ para todo $ n $, y
\emph{geométrico}\index{geométrico} con proporción común $ q $ si
$ u_{n+1} = q\,u_n $ para todo $ n $.
\end{definition}

\begin{proposition}[Forma explícita y sumas]\label{prop:g12:seq:arith-geom}
Sea $ n \in \N $.
\begin{enumerate}
  \item Si $(u_n)$ es \omterm{def:g12:seq:arith-geom}{aritmética} con \omterm{def:g11:seq:arithmetic}{diferencia común} $ r $, entonces
        $ u_n = u_0 + nr $ y
        \[
        u_0 + u_1 + \dots + u_n = (n+1)\,\frac{u_0 + u_n}{2}.
        \]
  \item Si $(u_n)$ es \omterm{def:g12:seq:arith-geom}{geométrico} con proporción común $ q \neq 1$, entonces
        $ u_n = u_0\, q^n $ y
        \[
        u_0 + u_1 + \dots + u_n = u_0\,\frac{1 - q^{n+1}}{1 - q}.
        \]
\end{enumerate}
\end{proposition}

\begin{proof}
Las formas explícitas siguen por inducciones inmediatas. Para el \omterm{def:g12:seq:arith-geom}{aritmética} suma,
escribir $ S = u_0 + \dots + u_n $ y sumar la misma suma escrita en orden inverso:
cada uno de los $ n+1$ sumas de columnas son iguales $ u_0 + u_n $, entonces $2S = (n+1)(u_0+u_n)$.
Para el \omterm{def:g12:seq:arith-geom}{geométrico} suma, calcular $ S - qS $: todos los términos se cancelan en pares excepto el
primero y el último, entonces $(1-q)S = u_0(1 - q^{n+1})$.
\end{proof}

\section{Límite de una sucesión}

\begin{definition}[sucesión convergente]\label{def:g12:seq:limit}
A \omterm{def:g12:seq:sequence}{sucesión} $(u_n)$ \emph{converger}\index{converger} hacia real
numero $\ell $ si cada abierto \omterm{def:g10:numbers:interval}{intervalo} que contiene $\ell $ contiene todos los términos
$ u_n $ a partir de algún índice en adelante. luego escribimos $\lim\limits_{n\to+\infty} u_n = \ell $.

Equivalentemente: por cada $\varepsilon > 0$, existe $ N \in \N $ tal que
para todos $ n \geq N $, $\abs{u_n - \ell} \leq \varepsilon $.
\end{definition}

\begin{omfigure}
\begin{tikzpicture}
\begin{axis}[omaxis,
  width=11cm, height=5.6cm,
  xlabel={$ n $}, ylabel={$ u_n $},
  xmin=0, xmax=16.5, ymin=0.7, ymax=2.9,
  xtick={4}, xticklabels={$ N $}, ytick={2}, yticklabels={$\ell $}]
  \fill[omDef!20] (axis cs:0,1.75) rectangle (axis cs:16.5,2.25);
  \addplot[omProp, dashed, domain=0:16.5] {2};
  \addplot[omDef, only marks, mark=*, mark size=1.5pt,
    samples at={1,2,3,4,5,6,7,8,9,10,11,12,13,14,15,16}] {2 + (-1)^x/x};
  \draw[black, dashed, thin] (axis cs:4,0.7) -- (axis cs:4,2.25);
  \node[font=\small, anchor=east] at (axis cs:16.3,2.4) {$\ell+\varepsilon $};
  \node[font=\small, anchor=east] at (axis cs:16.3,1.6) {$\ell-\varepsilon $};
\end{axis}
\end{tikzpicture}

{\small Convergencia de $ u_n = 2 + \frac{(-1)^n}{n}$ a $\ell = 2$: dado
$\varepsilon > 0$, todos los términos del índice $ N $ en mentira en la banda
$\intcc{\ell-\varepsilon}{\ell+\varepsilon}$.}
\end{omfigure}

\begin{definition}[Divergencia al infinito]\label{def:g12:seq:infinity}
La \omterm{def:g12:seq:sequence}{sucesión} $(u_n)$ \emph{tiende a $+\infty $} si por cada $ A \in \R $, allí
existe $ N \in \N $ tal que $ u_n \geq A $ para todo $ n \geq N $. escribimos
$\lim\limits_{n\to+\infty} u_n = +\infty $; la definición de
$\lim u_n = -\infty $ es analogo. A \omterm{def:g12:seq:sequence}{sucesión} que no converge se dice
a \emph{divergir}.
\end{definition}

\begin{remark}
A \omterm{def:g12:seq:sequence}{sucesión} puede divergir sin tender a $\pm\infty $: la \omterm{def:g12:seq:sequence}{sucesión}
$ u_n = (-1)^n $ toma solo los valores $1$ y $-1$ y no tiene límite.
\end{remark}

\begin{proposition}[Unicidad del límite]\label{prop:g12:seq:uniqueness}
Si $(u_n)$ \omterm{def:g12:seq:limit}{converger}, su límite es único.
\end{proposition}

\begin{proof}
Suponer $ u_n \to \ell $ y $ u_n \to \ell'$ con $\ell \neq \ell'$, decir
$\ell < \ell'$. Colocar $\varepsilon = \frac{\ell' - \ell}{3} > 0$. De algunos
índice encendido, $\abs{u_n - \ell} \leq \varepsilon $ y
$\abs{u_n - \ell'} \leq \varepsilon $, por eso
\[
\ell' - \ell \leq \abs{\ell' - u_n} + \abs{u_n - \ell}
\leq 2\varepsilon = \tfrac{2}{3}(\ell' - \ell) < \ell' - \ell,
\]
una contradicción.
\end{proof}

\begin{proposition}[Operaciones sobre límites]\label{prop:g12:seq:operations}
Sea $(u_n)$ y $(v_n)$ ser sucesiones con limites $\ell $ y $\ell'$ (finito o
infinito). Entonces, siempre que el lado derecho no sea una forma indeterminada,
\[
\lim (u_n + v_n) = \ell + \ell', \qquad
\lim (u_n v_n) = \ell\,\ell', \qquad
\lim \frac{u_n}{v_n} = \frac{\ell}{\ell'}.
\]
Las formas indeterminadas son
$(+\infty) + (-\infty)$, $0 \times \infty $, $\frac{\infty}{\infty}$ y
$\frac{0}{0}$.
\end{proposition}

\begin{proof}
Probamos la regla de la suma para límites finitos; los otros casos son similares y quedan
como ejercicios. Sea $\varepsilon > 0$. existen $ N_1, N_2$ tal que
$\abs{u_n - \ell} \leq \varepsilon/2$ para $ n \geq N_1$ y
$\abs{v_n - \ell'} \leq \varepsilon/2$ para $ n \geq N_2$. Para
$ n \geq \max(N_1, N_2)$, la desigualdad del triángulo da
\[
\abs{(u_n + v_n) - (\ell + \ell')}
\leq \abs{u_n - \ell} + \abs{v_n - \ell'} \leq \varepsilon. \qedhere
\]
\end{proof}

\begin{method}[Levantando una forma indeterminada]\label{met:g12:seq:indeterminate}
Ante una forma indeterminada, factorice el término dominante. Por ejemplo
\[
n^2 - n = n^2\left(1 - \tfrac{1}{n}\right) \xrightarrow[n\to+\infty]{} +\infty,
\qquad
\frac{2n^2+1}{n^2 - n} = \frac{2 + 1/n^2}{1 - 1/n}
\xrightarrow[n\to+\infty]{} 2 .
\]
\end{method}

\section{Teoremas de convergencia}

\begin{theorem}[Teoremas de comparación y compresión]\label{thm:g12:seq:squeeze}
Sea $(u_n)$, $(v_n)$, $(w_n)$ ser sucesiones.
\begin{enumerate}
  \item Si $ u_n \leq v_n $ de algún índice en adelante y $ u_n \to +\infty $, entonces
        $ v_n \to +\infty $.
  \item \emph{(Teorema de compresión)}\index{(Teorema de compresión)} Si
        $ u_n \leq v_n \leq w_n $ de algún índice en adelante y $(u_n)$ y $(w_n)$
        ambos convergen al mismo límite $\ell $, entonces $(v_n)$ \omterm{def:g12:seq:limit}{converger}
        a ~$\ell $.
\end{enumerate}
\end{theorem}

\begin{proof}
\emph{1.} Sea $ A \in \R $. Desde $ u_n \to +\infty $, hay $ N $ con
$ u_n \geq A $ para $ n \geq N $; ampliando $ N $ si es necesario, $ v_n \geq u_n \geq A $
para $ n \geq N $.

\emph{2.} Sea $\varepsilon > 0$. A partir de algún índice, ambos
$\ell - \varepsilon \leq u_n $ y $ w_n \leq \ell + \varepsilon $ espera, por lo tanto
$\ell - \varepsilon \leq u_n \leq v_n \leq w_n \leq \ell + \varepsilon $, eso es $\abs{v_n - \ell} \leq \varepsilon $.
\end{proof}

\begin{example}
A pesar de $ n \geq 1$, $-\frac{1}{n} \leq \frac{(-1)^n}{n} \leq \frac{1}{n}$, y
ambos límites tienden a $0$; por eso $\frac{(-1)^n}{n} \to 0$.
\end{example}

\begin{theorem}[Teorema de convergencia monótona]\label{thm:g12:seq:monotone}
Un \omterm{def:g12:seq:monotonic}{creciente} \omterm{def:g12:seq:sequence}{sucesión} eso es \omterm{def:g12:seq:bounded}{delimitado arriba} \omterm{def:g12:seq:limit}{converger}. A \omterm{def:g10:functions:variations}{decreciente} \omterm{def:g12:seq:sequence}{sucesión}
eso es delimitado por debajo \omterm{def:g12:seq:limit}{converger}. Un \omterm{def:g12:seq:monotonic}{creciente} \omterm{def:g12:seq:sequence}{sucesión} eso no es \omterm{def:g12:seq:bounded}{delimitado
arriba} tiende a $+\infty $.
\end{theorem}

\begin{proof}[Prueba parcial]
Probamos la tercera afirmación. Sea $(u_n)$ ser \omterm{def:g12:seq:monotonic}{creciente} y no \omterm{def:g12:seq:bounded}{delimitado
arriba}, y deja $ A \in \R $. Desde $ A $ no hay un límite superior, existe $ N $
con $ u_N \geq A $; por monotonicidad, $ u_n \geq u_N \geq A $ para todo $ n \geq N $.
Por eso $ u_n \to +\infty $.

Las dos declaraciones de convergencia se basan en la propiedad del límite superior \omterm{def:g10:functions:extrema}{mínimo} de
$\R $; ellos son \emph{admitido en este nivel} (y demostrado en el primer año de
universidad).
\end{proof}

\begin{remark}
El teorema garantiza la \emph{existencia} del límite pero no da
su valor. Un \omterm{def:g12:seq:monotonic}{creciente} \omterm{def:g12:seq:sequence}{sucesión} \omterm{def:g12:seq:bounded}{delimitado arriba} por $ M $ \omterm{def:g12:seq:limit}{converger} a algunos
$\ell \leq M $, no necesariamente a $ M $.
\end{remark}

\begin{theorem}[Límite de sucesiones geométricas.]\label{thm:g12:seq:geometric}
Sea $ q \in \R $.
\begin{enumerate}
  \item Si $ q > 1$, entonces $ q^n \to +\infty $.
  \item Si $ q = 1$, entonces $ q^n \to 1$.
  \item Si $\abs{q} < 1$, entonces $ q^n \to 0$.
  \item Si $ q \leq -1$, entonces $(q^n)$ diverge y no tiene límite.
\end{enumerate}
\end{theorem}

\begin{proof}
\emph{1.} Escribir $ q = 1 + a $ con $ a > 0$. La desigualdad de Bernoulli
(\cref{ex:g12:seq:bernoulli}) da $ q^n \geq 1 + na \to +\infty $, y nosotros
concluir por comparación (\cref{thm:g12:seq:squeeze}).

\emph{2.} Inmediato.

\emph{3.} Si $ q = 0$ el reclamo es claro. De lo contrario $\abs{q} < 1$ da
$1/\abs{q} > 1$, entonces $(1/\abs{q})^n \to +\infty $ por el punto 1, por lo tanto
$\abs{q}^n \to 0$, y $-\abs{q}^n \leq q^n \leq \abs{q}^n $ nos permite
concluir por el teorema de compresión.

\emph{4.} Para $ q \leq -1$, $(q^{2n})$ toma valores $\geq 1$ mientras
$(q^{2n+1})$ toma valores $\leq -1$: ningún límite puede atraer a ambos
subsecuencias.
\end{proof}

\begin{omfigure}
\begin{tikzpicture}
\begin{axis}[omaxis,
  width=10cm, height=5.6cm,
  xlabel={$ n $}, ylabel={$ q^n $},
  xmin=-0.5, xmax=13, ymin=-1.3, ymax=4.3,
  xtick={5,10}, ytick={-1,1}]
  \addplot[omThm, only marks, mark=*, mark size=1.5pt,
    samples at={0,1,2,3,4,5,6}] {1.25^x};
  \addplot[omDef, only marks, mark=*, mark size=1.5pt,
    samples at={0,1,2,3,4,5,6,7,8,9,10,11,12}] {0.75^x};
  \addplot[omProp, only marks, mark=*, mark size=1.5pt,
    samples at={0,1,2,3,4,5,6,7,8,9,10,11,12}] {(-0.8)^x};
  \node[omThm, font=\small, anchor=west] at (axis cs:6.3,3.7) {$ q=1.25$};
  \node[omDef, font=\small, anchor=south west] at (axis cs:3.4,0.55) {$ q=0.75$};
  \node[omProp, font=\small, anchor=north] at (axis cs:7,-0.75) {$ q=-0.8$};
\end{axis}
\end{tikzpicture}

{\small Los tres comportamientos de $(q^n)$: divergencia a $+\infty $ para
$ q > 1$ (rojo), convergencia a $0$ para $\abs q < 1$ (azul) y amortiguado
oscilación --- todavía convergencia a $0$ --- para $-1 < q < 0$ (naranja).}
\end{omfigure}

\begin{method}[Sucesiones recurrentes $ u_{n+1} = f(u_n)$]\label{met:g12:seq:recurrent}
para estudiar un \omterm{def:g12:seq:sequence}{sucesión} definido por $ u_{n+1} = f(u_n)$:
\begin{enumerate}
  \item probar por inducción que $(u_n)$ se queda en un \omterm{def:g10:numbers:interval}{intervalo} $ I $ en el cual $ f $
        se porta bien (y, a menudo, eso $(u_n)$ es monótona);
  \item deducir la convergencia del teorema de convergencia monótona;
  \item pasar al límite en la relación $ u_{n+1} = f(u_n)$: si $ f $ es
        continuo y $ u_n \to \ell \in I $, entonces $\ell $ satisface
        $ f(\ell) = \ell $ (ver \cref{ch:g12:limcont}); resuelve esto \omterm{def:g10:algebra:equation}{ecuación} y
        seleccione el derecho \omterm{def:g11:quad:discriminant}{raíz}.
\end{enumerate}
\end{method}

\begin{omfigure}
\begin{tikzpicture}
\begin{axis}[omaxis,
  width=8.4cm, height=6.4cm,
  xmin=-0.4, xmax=3.1, ymin=-0.4, ymax=2.7,
  xtick={0,1.414,1.848}, xticklabels={$ u_0$,$ u_1$,$ u_2$}, ytick=\empty]
  \addplot[omProp, dashed, domain=-0.3:2.6] {x};
  \addplot[omDef, thick, domain=-0.4:2.9] {sqrt(x+2)};
  \draw[omThm, thick] (axis cs:0,0) -- (axis cs:0,1.414)
    -- (axis cs:1.414,1.414) -- (axis cs:1.414,1.848)
    -- (axis cs:1.848,1.848) -- (axis cs:1.848,1.961)
    -- (axis cs:1.961,1.961);
  \draw[black, dashed, thin] (axis cs:1.414,0) -- (axis cs:1.414,1.414);
  \draw[black, dashed, thin] (axis cs:1.848,0) -- (axis cs:1.848,1.848);
  \fill (axis cs:2,2) circle (1.5pt)
    node[below right, font=\small] {$\ell = 2$};
  \node[omDef, font=\small, above left] at (axis cs:2.75,2.18)
    {$ y = f(x)$};
  \node[omProp, font=\small, anchor=north west] at (axis cs:1.08,0.82)
    {$ y = x $};
\end{axis}
\end{tikzpicture}

{\small Construcción de escaleras para $ u_{n+1} = \sqrt{u_n + 2}$, $ u_0 = 0$
(\cref{exo:g12:seq:6}): cada paso vertical lee $ f(u_n)$ en la curva,
cada paso horizontal lo lleva de regreso $ y = x $. La \omterm{def:g12:seq:sequence}{sucesión} sube
al \omterm{pb:g10:functions:1}{punto fijo} $\ell = 2$, donde la curva se encuentra con la línea.}
\end{omfigure}

\begin{example}\label{ex:g12:seq:sqrt2}
Sea $ u_0 = 2$ y $ u_{n+1} = \frac{1}{2}\left(u_n + \frac{2}{u_n}\right)$.
Se comprueba por inducción que $ u_n \geq \sqrt{2}$ para todo $ n $ (la desigualdad
$\frac{1}{2}(x + 2/x) \geq \sqrt{2}$ para $ x>0$ es equivalente a
$(x - \sqrt2)^2 \geq 0$), entonces eso $(u_n)$ es \omterm{def:g10:functions:variations}{decreciente}, desde
\[
u_{n+1}-u_n=\frac{2-u_n^2}{2u_n}\leq 0 .
\]
\omterm{def:g10:functions:variations}{decreciente} y delimitado por debajo, $(u_n)$
\omterm{def:g12:seq:limit}{converger} a algunos $\ell \geq \sqrt{2}$, que debe satisfacer
$\ell = \frac{1}{2}(\ell + 2/\ell)$, es decir.\ $\ell^2 = 2$. Por eso
$ u_n \to \sqrt{2}$. Este es el algoritmo de Heron, ya utilizado por el
babilonios; su convergencia es extremadamente rápida ($ u_3$ ya da $\sqrt 2$
hasta ocho decimales).
\end{example}

\section{Ejercicios}

\begin{exercise}[$\star $]\label{exo:g12:seq:1}
Demostrar por inducción que para todos $ n \in \N $,
\[
1^2 + 2^2 + \dots + n^2 = \frac{n(n+1)(2n+1)}{6}.
\]
\end{exercise}

\begin{exercise}[$\star $]\label{exo:g12:seq:2}
Estudiar la monotonicidad de la sucesiones definido para $ n \geq 1$ por
\[
a_n = \frac{n+1}{n}, \qquad
b_n = \frac{2^n}{n}, \qquad
c_n = n^2 - 10n .
\]
\end{exercise}

\begin{exercise}[$\star $]\label{exo:g12:seq:3}
Calcule los límites de la sucesiones con términos generales
\[
u_n = \frac{3n^2 - n + 1}{2n^2 + 5}, \qquad
v_n = \sqrt{n+1} - \sqrt{n}, \qquad
w_n = \frac{2^n - 3^n}{3^n + 1}.
\]
\end{exercise}

\begin{exercise}[$\star $]\label{exo:g12:seq:4}
Sea $(u_n)$ ser el \omterm{def:g12:seq:sequence}{sucesión} \omterm{def:g12:seq:arith-geom}{aritmética} con $ u_0 = 5$ y \omterm{def:g11:seq:arithmetic}{diferencia común}
$ r = 3$, y $(v_n)$ el \omterm{def:g12:seq:sequence}{sucesión} geométrica con $ v_0 = 8$ y proporción común
$ q = \frac{1}{2}$. Calcular $ u_n $, $ v_n $, $\sum_{k=0}^{n} u_k $ y
$\sum_{k=0}^{n} v_k $, y los límites de las cuatro expresiones como
$ n \to +\infty $.
\end{exercise}

\begin{exercise}[$\star\star $]\label{exo:g12:seq:5}
Utilizando el teorema de compresión, calcule
\[
\lim_{n\to+\infty} \frac{n + \cos n}{n + 1}
\qquad\text{and}\qquad
\lim_{n\to+\infty} \frac{n!}{n^n},
\]
dónde $ n! = 1 \times 2 \times \dots \times n $. Para el segundo límite, obligado
$\frac{n!}{n^n}$ por un término de un \omterm{def:g12:seq:sequence}{sucesión} geométrica.
\end{exercise}

\begin{exercise}[$\star\star $]\label{exo:g12:seq:6}
Sea $ u_0 = 0$ y $ u_{n+1} = \sqrt{u_n + 2}$ para todo $ n \in \N $.
\begin{enumerate}
  \item Demostrar por inducción que $0 \leq u_n \leq 2$ para todo $ n $.
  \item mostrar que $(u_n)$ es \omterm{def:g12:seq:monotonic}{creciente}.
  \item deducir eso $(u_n)$ \omterm{def:g12:seq:limit}{converger} y determinar su límite.
\end{enumerate}
\end{exercise}

\begin{exercise}[$\star\star $]\label{exo:g12:seq:7}
Un paciente toma una dosis de $1$ unidad de medicamento cada mañana. Durante cada
En un periodo de 24 horas, el cuerpo elimina $40\%$ del fármaco presente. Sea $ u_n $ ser
la cantidad de fármaco en el cuerpo justo después de la dosis del día $ n $, de modo que
$ u_0 = 1$.
\begin{enumerate}
  \item justifica eso $ u_{n+1} = 0.6\,u_n + 1$.
  \item Sea $ v_n = u_n - 2.5$. mostrar que $(v_n)$ es \omterm{def:g12:seq:arith-geom}{geométrico} y deducir un
        fórmula explícita para $ u_n $.
  \item Determinar la cantidad a largo plazo de fármaco en el cuerpo.
\end{enumerate}
\end{exercise}

\begin{exercise}[$\star\star $]\label{exo:g12:seq:8}
Sea $(u_n)$ ser definido por $ u_0 = 3$ y $ u_{n+1} = \frac{4u_n - 1}{u_n + 2}$.
\begin{enumerate}
  \item Demuestre por inducción que $ u_n > 1$ para todo $ n \in \N $.
  \item mostrar que $ v_n = \dfrac{1}{u_n - 1}$ define un \omterm{def:g12:seq:sequence}{sucesión} \omterm{def:g12:seq:arith-geom}{aritmética}.
  \item Deducir fórmulas explícitas para $ v_n $ y $ u_n $, y el límite de
        $(u_n)$.
\end{enumerate}
\end{exercise}

\begin{exercise}[$\star\star\star $]\label{exo:g12:seq:9}
Para $ n \geq 1$, dejar
$ H_n = 1 + \frac{1}{2} + \frac{1}{3} + \dots + \frac{1}{n}$.
\begin{enumerate}
  \item \omterm{def:g10:proba:sample}{Muestra} eso para todos $ n \geq 1$, $ H_{2n} - H_n \geq \frac{1}{2}$.
  \item deducir eso $ H_{2^k} \geq 1 + \frac{k}{2}$ para todo $ k \in \N $, y
        concluir que $ H_n \to +\infty $.
\end{enumerate}
\end{exercise}

\begin{exercise}[$\star\star\star $]\label{exo:g12:seq:10}
\emph{(Adyacente sucesiones.)} Dos sucesiones $(a_n)$ y $(b_n)$ son
\emph{adyacente} si $(a_n)$ es \omterm{def:g12:seq:monotonic}{creciente}, $(b_n)$ es \omterm{def:g10:functions:variations}{decreciente}, y
$ b_n - a_n \to 0$.
\begin{enumerate}
  \item \omterm{def:g10:proba:sample}{Muestra} eso para todos $ n $, $ a_n \leq b_n $. (Pista: estudia la monotonicidad
        de $(b_n - a_n)$.)
  \item Muestre que adyacente sucesiones ambos convergen, al mismo límite.
  \item Aplicación: demostrar que el sucesiones
        $ a_n = \sum_{k=0}^{n} \frac{1}{k!}$ y
        $ b_n = a_n + \frac{1}{n \cdot n!}$ ($ n \geq 1$) son adyacentes. (Sus
        El límite común es el número. $\eu $, estudiado en \cref{ch:g12:exp}.)
\end{enumerate}
\end{exercise}

\section{Problema: la sucesión de Heron, finalmente juzgada}

\begin{problem}[{Problema del fin de semana --- certificación de inducción, monótono
sentencias de convergencia y la receta de hace dos mil años para
$\sqrt2$ finalmente obtiene su prueba (con el maravilloso medio de Gauss para
postre)}]
\label{pb:g12:seq:1}
Tres veces esta serie ha cumplido con la receta de Heron --- promedia el
adivina con $2/\text{guess}$ --- y tres veces solo pudo
\emph{observar} que la receta funciona. Este capítulo por fin posee
los instrumentos de juicio: la inducción
(\cref{thm:g12:seq:induction}), el teorema de convergencia monótona
(\cref{thm:g12:seq:monotone}) y límites de recurrencias. el
veredicto y la velocidad certificada ocupan el centro de este
problema; A su alrededor, las clásicas trampas de la inducción, las más lentas.
divergencia en matemáticas, y la convergencia más rápida de Gauss jamás
encontrado.

\medskip
\noindent\textbf{Part I --- Induction warm-ups.}
\begin{enumerate}
  \item Demostrar por inducción:
        $1 + 3 + 5 + \dots + (2n - 1) = n^2$ (la escalera de
        números \omterm{def:g11:func:parity}{impares}, dibujados en el volumen de la Escuela Secundaria, ahora
        certificado).
  \item Demostrar por inducción que $2^n > n $ por cada
        $ n \in \N $.
  \item Demuestre la desigualdad de Bernoulli por inducción: para
        $ x \geq 0$ y $ n \in \N $,
        $(1 + x)^n \geq 1 + nx $.
  \item La trampa clásica: ``todas las canicas tienen el mismo color ---
        cierto para una canica; y si alguno $ n $ las canicas son siempre
        monocromo, luego entre $ n + 1$ canicas la primera $ n $
        compartir un color, el último $ n $ comparten un color, así que todos
        $ n + 1$ hacer". Cada niño sabe que la conclusión es
        Absurdo: encontrar el paso exacto donde se rompe la inducción.
  \item Demostrar por inducción que $4^n - 1$ es divisible por $3$
        por cada $ n \in \N $.
\end{enumerate}

\medskip
\noindent\textbf{Part II --- The trial of Heron.}
Sea $ x_0 = 2$ y $ x_{n+1} = \dfrac12\left(x_n +
\dfrac{2}{x_n}\right)$.
\begin{enumerate}[resume]
  \item Calcular $ x_1$, $ x_2$, $ x_3$ como fracciones exactas (antiguo
        amigos).
  \item Demostrar la identidad clave
        \[
        x_{n+1}^2 - 2 = \left(\frac{x_n^2 - 2}{2x_n}\right)^{\!2}
        \geq 0,
        \]
        y deducir por inducción que $ x_n > 0$ y
        $ x_n^2 > 2$ por cada $ n $.
  \item mostrar que $(x_n)$ es estrictamente \omterm{def:g10:functions:variations}{decreciente} (calcular
        $ x_{n+1} - x_n $ y utilice la pregunta 7).
  \item Invoque el teorema de convergencia monótona: ¿por qué
        $(x_n)$ \omterm{def:g12:seq:limit}{converger} a algún límite $ L \geq 1$?
  \item Identificar el límite: pasar la recurrencia al límite
        (\cref{prop:g12:seq:operations}) y concluir
        $ L = \sqrt2$. Enuncia el veredicto histórico: después de dos
        mil años de fiel servicio, la receta de Heron es
        \emph{probado} \omterm{def:g12:seq:limit}{converger}.
  \item La velocidad certificada: con $ e_n = x_n - \sqrt2$, probar
        \[
        e_{n+1} = \frac{e_n^2}{2 x_n} ,
        \]
        y deducir $ e_{n+1} \leq \frac{e_n^2}{2\sqrt2}$: el
        el error se eleva al cuadrado en cada paso --- la duplicación de dígitos
        observado desde el volumen de la escuela secundaria, ahora un teorema.
  \item Confirmar numéricamente: calcular $ e_0, e_1, e_2, e_3$ (de
        pregunta 6) y comprobar que cada $\frac{e_{n+1}}{e_n^2}$
        esta cerca de $\frac{1}{2x_n}$.
\end{enumerate}

\medskip
\noindent\textbf{Part III --- The slowest divergence.}
\begin{enumerate}[resume]
  \item \Cref{exo:g12:seq:9} probado
        $ H_{2^k} \geq 1 + \frac k2$ para las sumas armónicas. como
        muchos términos garantizan $ H_n > 10$? (Una potencia de dos voluntades
        hacer; Maravíllate con su tamaño.)
  \item Por el contrario, el \omterm{def:g12:seq:arith-geom}{geométrico} sumas
        $1 + \frac12 + \frac14 + \dots + \frac{1}{2^n} =
        2 - \frac{1}{2^n}$ \omterm{def:g12:seq:limit}{converger} a $2$
        (\cref{thm:g12:seq:geometric}): la barra de chocolate
        Intuición del volumen de la Escuela Secundaria, finalmente un límite.
        declaración. Escribe la prueba de dos líneas.
  \item Entre los dos: demuestra que las sumas
        $ S_n = 1 + \frac{1}{4} + \frac{1}{9} + \dots +
        \frac{1}{n^2}$ \omterm{def:g12:seq:limit}{converger}, limitando
        $\frac{1}{k^2} \leq \frac{1}{k(k-1)} =
        \frac{1}{k-1} - \frac{1}{k}$ (para $ k \geq 2$),
        telescópico y aplicación de convergencia monótona. (El
        límite, $\frac{\pi^2}{6}$, es uno de los milagros de Euler,
        demostrado en los volúmenes de la universidad.)
  \item Indique la moraleja de las preguntas 13 a 15 en dos oraciones:
        ¿A qué tienden los términos? $0$'' decidir sobre el
        convergencia de las sumas --- ¿y qué no?
\end{enumerate}

\medskip
\noindent\textbf{Part IV --- Gauss's arithmetic--geometric
\omterm{def:g11:stat:mean}{media}.}
Sea $ a_0 = 1$, $ b_0 = 2$, y
\[
a_{n+1} = \sqrt{a_n b_n},
\qquad
b_{n+1} = \frac{a_n + b_n}{2} .
\]
\begin{enumerate}[resume]
  \item Calcular $ a_1, b_1, a_2, b_2$ (cinco decimales). ¿Qué hacer?
        observas sobre la velocidad?
  \item mostrar que $ a_n \leq b_n $ por cada $ n $ (el
        aritmética--desigualdad geométrica, encontrada a lo largo de este
        serie), que $(a_n)$ aumenta y $(b_n)$ disminuye.
  \item mostrar que
        $ b_{n+1} - a_{n+1} \leq \frac{b_n - a_n}{2}$
        (factor $ b_{n+1} - a_{n+1} =
        \frac{(\sqrt{b_n} - \sqrt{a_n})^2}{2}$ y comparar),
        y concluir con \cref{exo:g12:seq:10} que los dos
        sucesiones son adyacentes: comparten un límite común
        $ M(1, 2)$, el \emph{aritmética--geométrica \omterm{def:g11:stat:mean}{media}}.
  \item Calcular $ M(1, 2)$ a seis decimales (cuántas iteraciones
        necesitabas?). El 30 de mayo de 1799, Gauss calculó
        $ M(1, \sqrt2)$ hasta once decimales, reconocido
        $\frac{\pi}{M(1,\sqrt2)}$ como una integral conocida, y
        escribió que se había abierto un "nuevo campo de análisis" ---
        tenía: integrales elípticas, contadas en la universidad
        volúmenes. Cierre con las velocidades de convergencia observadas de
        este problema, del más lento al más rápido.
\end{enumerate}
\end{problem}
